% The two rule graphs in Lemma 3.17. Red marks the substituted edges/copies.
\begin{tikzpicture}[
  x=1cm,y=1cm,line cap=round,line join=round,
  edge/.style={line width=.75pt},
  replaced/.style={draw=red!45!black,line width=.9pt},
  vertex/.style={circle,fill=black,inner sep=1.65pt},
  terminal/.style={circle,fill=black,inner sep=2.5pt},
  newvertex/.style={circle,fill=red!45!black,inner sep=1.65pt},
  every node/.style={font=\small}
]
% A = W: highlight the opposite edges su and vt.
\begin{scope}
  \coordinate (s) at (-1.8,0);
  \coordinate (u) at (0,1.4);
  \coordinate (t) at (1.8,0);
  \coordinate (v) at (0,-1.4);
  \draw[edge] (u)--(t) (s)--(v) (u)--(v);
  \draw[replaced] (s)--(u) (v)--(t);
  \node[terminal,label=left:$s$] at (s) {};
  \node[terminal,label=right:$t$] at (t) {};
  \node[vertex,label=above:$u$] at (u) {};
  \node[vertex,label=below:$v$] at (v) {};
  \node at (0,-2.05) {$A=W$};
\end{scope}
% Replace each of the two marked edges by a fresh Wheatstone graph.
\draw[-{Stealth[length=2mm,width=1.6mm]},black!55,line width=.65pt]
  (2.45,0)--node[above=5pt,text=black]{$su,vt\mapsto W$}(4.35,0);
% B: the two red copies have disjoint new internal vertices.
\begin{scope}[xshift=6.8cm]
  \coordinate (s) at (-1.8,0);
  \coordinate (u) at (0,1.4);
  \coordinate (t) at (1.8,0);
  \coordinate (v) at (0,-1.4);
  \coordinate (su1) at (-1.25,1.15);
  \coordinate (su2) at (-.55,.25);
  \coordinate (vt1) at (.55,-.25);
  \coordinate (vt2) at (1.25,-1.15);
  \draw[edge] (u)--(t) (s)--(v) (u)--(v);
  \draw[replaced] (s)--(su1)--(u)--(su2)--(s) (su1)--(su2);
  \draw[replaced] (v)--(vt1)--(t)--(vt2)--(v) (vt1)--(vt2);
  \foreach \p in {su1,su2,vt1,vt2}
    \node[newvertex] at (\p) {};
  \node[terminal,label=left:$s$] at (s) {};
  \node[terminal,label=right:$t$] at (t) {};
  \node[vertex,label=above:$u$] at (u) {};
  \node[vertex,label=below:$v$] at (v) {};
  \node at (0,-2.05) {$B$};
\end{scope}
\end{tikzpicture}
